%! Dimensions of the Four Subspaces — Unit 3, Lesson 3.5 — Warm-Up (Answer Key)
\documentclass[10pt]{article}
\usepackage{linalg-article}
\usepackage{linalg-key}   % pulls in -boxes; do NOT also load -boxes
\newcommand{\TallMath}[1]{$\displaystyle #1\rule[-1.4em]{0pt}{3.2em}$}
\newcommand{\xx}{\mathbf{x}}
\newcommand{\svec}{\mathbf{s}}
\newcommand{\zero}{\mathbf{0}}
\newcommand{\T}{\mathsf{T}}

\begin{document}

\pageheader{Unit 3, Lesson 3.5}{Warm-Up --- Answer Key}
\namedateperiod

\vspace{0.12in}

\begin{notesbox}{Getting Ready: Skills We'll Use Today}
\small These skills from Lessons 2.4 and 3.4 power today's work. Answer quickly.

\vspace{4pt}
\textbf{1. Transpose (Lesson 2.4).} The transpose $A^{\T}$ turns the rows of $A$ into columns. Fill it in.
\[
  A=\begin{bmatrix} 1 & 1 & 2 \\ 2 & 1 & 3 \end{bmatrix}
  \qquad
  A^{\T}=\begin{bmatrix}
    \rule{0pt}{1.1em}{\color{keyred}\mathbf{1}} & {\color{keyred}\mathbf{2}} \\[3pt]
    {\color{keyred}\mathbf{1}} & {\color{keyred}\mathbf{1}} \\[3pt]
    {\color{keyred}\mathbf{2}} & {\color{keyred}\mathbf{3}}
  \end{bmatrix}
\]

\vspace{4pt}
\textbf{2. Rank and shape (Lesson 3.4).} The matrix below is already reduced. Give the number of rows $m$,
columns $n$, and pivots $r$.
\[
  \begin{bmatrix} 1 & 0 & 1 \\ 0 & 1 & 1 \\ 0 & 0 & 0 \end{bmatrix}
  \qquad
  m=\ans{$3$}\quad n=\ans{$3$}\quad r=\ans{$2$}
\]

\vspace{4pt}
\textbf{3. Nullspace dimension (Lesson 3.4).} Using the numbers from item~2, the nullspace of that matrix
has dimension $\dim N(A)=n-r=\ans{$1$}$. How many free columns is that? \ans{$1$ (column 3)}
\end{notesbox}

\begin{teachernote}
Each item is a number today's lesson needs. Item~1 previews how the \emph{row} space and \emph{left}
nullspace are really ``$C(A^{\T})$ and $N(A^{\T})$'' --- transpose, then use the tools you already have.
Item~2's three numbers $m=3$, $n=3$, $r=2$ are exactly the ingredients of every dimension: $\dim C(A)=\dim
C(A^{\T})=r=2$, $\dim N(A)=n-r=1$, $\dim N(A^{\T})=m-r=1$. Item~3 is the nullspace count they already know
($n-r$); today's new count is the \emph{left} nullspace $m-r$.
\end{teachernote}

\end{document}
